% Part: lambda-calculus
% Chapter: lambda-definability
% Section: fixpoints

\documentclass[../../../include/open-logic-section]{subfiles}

\begin{document}

\olfileid[pt]{lam}{ldf}{fp}

\olsection{Pontos fixos}

Suponha que quiséssemos definir recursivamente a função fatorial como um
termo $\fn{Fac}$ com a seguinte propriedade:
\[
\fn{Fac} \ident \lambd[n][\fn{IsZero}\, n\, \num 1 (\fn{Mult}\, n (\fn{Fac}(\fn{Pred} \, n)))]
\]
Isto é, o fatorial de $n$ é $1$ se $n = 0$, e, caso contrário, é $n$ vezes o
fatorial de $n-1$. Naturalmente, não podemos definir o termo
$\fn{Fac}$ desse modo, pois o próprio $\fn{Fac}$ ocorre no lado
direito. Definições recursivas desse tipo, que envolvem autorreferência,
não fazem parte do cálculo lambda. Definir um termo, p.ex., por
\[
\fn{Mult} \ident \lambd[ab][a (\fn{Add}\, b) 0]
\]
envolve apenas termos previamente definidos no lado direito, como
$\fn{Add}$. Sempre podemos eliminar $\fn{Add}$ substituindo-o por seu
termo definidor. Isso daria o termo $\fn{Mult}$ como um termo lambda
puro. Se o próprio $\fn{Add}$ envolvesse termos definidos (como ocorre, p.ex.,
com $\fn{Add}'$), poderíamos continuar esse processo e finalmente chegar a
um termo lambda puro.

No entanto, isso não vale para definições recursivas como a de
$\fn{Fac}$ acima. Se substituirmos a ocorrência de $\fn{Fac}$ no
lado direito pela própria definição de $\fn{Fac}$, obtemos:
\begin{align*}
  \fn{Fac} & \ident \lambd[n][\fn{IsZero}\, n\, \num{1}] \\
    & \qquad (\fn{Mult}\, n
  ((\lambd[n][\fn{IsZero} \, n \, \num 1\, (\fn{Mult}\, n\, (\fn{Fac}
    (\fn{Pred}\, n)))]) (\fn{Pred}\, n)))
\end{align*}
e ainda não eliminamos $\fn{Fac}$ do lado direito.
É evidente que, se repetirmos esse processo, a definição ficará cada vez
mais longa e o processo nunca resultará em um termo lambda puro. Portanto, esse
modo de definir o fatorial (ou, mais geralmente, funções recursivas)
não é viável.

A definição recursiva, contudo, nos diz algo: se $f$ fosse um
termo que representasse a função fatorial, então o termo
\[
\fn{Fac}' \ident \lambd[g][\lambd[n][\fn{IsZero} \, n \, \num 1\, (\fn{Mult} \, n\, (g (\fn{Pred} n)))]]
\]
aplicado ao termo $f$, O.s.\@, $\fn{Fac}'\,f$, também representaria a
função fatorial. Isto é, se considerarmos $\fn{Fac}'$ uma função
que recebe uma função e devolve uma função, o valor de
$\fn{Fac'}\, f$ é simplesmente~$f$, desde que $f$ seja o fatorial. Uma
função~$f$ com a propriedade de que $\fn{Fac}' \, f \equal[\beta] f$ é chamada
\emph{ponto fixo} de $\fn{Fac}'$. Assim, o fatorial é um ponto fixo de
$\fn{Fac}'$.

Há termos no cálculo lambda que calculam os pontos fixos de um
termo dado, e esses termos podem então ser usados para transformar um termo como
$\fn{Fac}'$ em uma definição do fatorial. 

\begin{defn}
  \ollabel{defn:Turing-Y}
  O \emph{combinador Y} é o termo:
  \[
  Y \ident (\lambd[ux][x(uux)])(\lambd[ux][x(uux)]).
  \]
\end{defn}

\begin{thm}
  $Y$ tem a propriedade de que $Yg \red g(Yg)$ para qualquer termo $g$. Assim,
  $Yg$ é sempre um ponto fixo de~$g$.
\end{thm}

\begin{proof}
  Abreviemos $(\lambd[ux][x(uux)])$ por~$U$, de modo que $Y \ident UU$. Então
  \begin{align*}
    Y g &\ident (\lambd[ux][x(uux)])U\,g \\
    &\red (\lambd[x][x(UUx)])g\\
    & \red g(UUg) \ident g(Yg).
  \end{align*}
  Como $g(Yg)$ e $Yg$ se reduzem ambos a $g(Yg)$, $g(Yg) \equal[\beta]
  Yg$; portanto, $Yg$ é um ponto fixo de~$g$.
\end{proof}

Naturalmente, como $Yg$ é um redex, a redução pode continuar indefinidamente:
\begin{align*}
  Y g &\red g (Y g) \\
    &\red g (g (Y g)) \\
    &\red g(g (g (Y g)))\\
    &\ldots
\end{align*}
Assim, podemos pensar em $Yg$ como $g$ aplicado a si mesmo infinitas
vezes. Se lhe aplicarmos~$g$ uma vez mais, por assim
dizer, não faremos nada adicional; $g$~aplicado a~$g$, aplicado
infinitas vezes a~$Yg$, ainda é $g$~aplicado a~$Yg$ infinitas
vezes.

Observe que a sequência acima de passos de redução~$\beta$ que começa
com~$Yg$ é infinita. Assim, se aplicarmos $Yg$ a algum termo, O.s.\@,
se considerarmos $(Yg)N$, esse termo também se reduzirá a infinitos
termos diferentes, a saber, $(g(Yg))N$, $(g(g(Yg)))N$, \dots. Ainda assim,
é possível que alguma \emph{outra} sequência de passos de redução
termine em uma forma normal.

Considere, por exemplo, o fatorial. Defina $\fn{Fac}$ como $Y\, \fn{Fac}'$
(O.s.\@, um ponto fixo de $\fn{Fac'}$). Então:
\begin{align*}
  \fn{Fac}\,\num 3
  &\red Y\, \fn{Fac}' \, \num 3 \\
  &\red \fn{Fac}' (Y \,\fn{Fac}') \, \num 3 \\
  &\ident (\lambd[x][\lambd[n][\fn{IsZero} \, n\, \num 1\,
      (\fn{Mult}\, n\, (x (\fn{Pred}\, n)))]]) \, \fn{Fac} \, \num 3\\
  & \red \fn{IsZero} \, \num 3\, \num 1\,
      (\fn{Mult}\, \num 3\, (\fn{Fac} (\fn{Pred}\, \num 3))) \\
  &\red  \fn{Mult}\, \num 3 \, (\fn{Fac} \, \num 2).
  \intertext{De modo semelhante,}
  \fn{Fac}\,\num 2
  &\red  \fn{Mult}\, \num 2 \, (\fn{Fac} \, \num 1) \\
  \fn{Fac}\,\num 1
  &\red  \fn{Mult}\, \num 1 \, (\fn{Fac} \, \num 0)
  \intertext{mas}
  \fn{Fac}\,\num 0
  &\red \fn{Fac}' (Y \,\fn{Fac}') \, \num 0 \\
  &\ident (\lambd[x][\lambd[n][\fn{IsZero} \, n\, \num 1\,
      (\fn{Mult}\, n\, (x (\fn{Pred}\, n)))]]) \, \fn{Fac} \, \num 0\\
  & \red \fn{IsZero} \, \num 0\, \num 1\,
      (\fn{Mult}\, \num 0\, (\fn{Fac} (\fn{Pred}\, \num 0))).\\
  &\red \num 1.
  \intertext{Portanto, em conjunto,}
  \fn{Fac}\,\num 3
  &\red  \fn{Mult}\, \num 3 \,
  (\fn{Mult} \, \num 2\, (\fn{Mult} \, \num 1\, \num 1)).
\end{align*}

O que vale para $\fn{Fac'}$ vale para qualquer definição recursiva. Suponha que
tenhamos uma equação recursiva
\begin{align*}
g\,x_1\dots x_n & \equal[\beta] N
\intertext{onde $N$ pode conter~$g$ e $x_1$, \dots,~$x_n$. Então sempre existe
um termo~$G$ dado por}
G & \ident (Y \lambd[g][\lambd[x_1\dots x_n][N]])\\
\intertext{tal que}
G\,x_1 \dots x_n & \equal[\beta] \Subst{N}{G}{g}.
\intertext{Pois, pelo teorema do ponto fixo,}
G \ident (Y \lambd[g][\lambd[x_1\dots x_n][N]]) & \red \lambd[g][\lambd[x_1\dots x_n][N]](Y \lambd[g][\lambd[x_1\dots x_n][N]])\\
& \ident (\lambd[g][\lambd[x_1\dots x_n][N]])\,G
\intertext{e, consequentemente,}
G\,x_1 \dots x_n 
& \red (\lambd[g][\lambd[x_1\dots x_n][N]])\,G\,x_1\dots x_n\\
& \red (\lambd[x_1\dots x_n][\Subst{N}{G}{g}])\,x_1\dots x_n\\
  & \red \Subst{N}{G}{g}.
\end{align*}

O combinador $Y$ de \olref{defn:Turing-Y} deve-se a Alan
Turing. Alonzo Church havia proposto uma versão diferente, que chamaremos
de~$Y_C$:
\[
Y_C \ident \lambd[g][(\lambd[x][g(xx)])(\lambd[x][g(xx)])].
\]
O combinador de Church é um pouco mais fraco que o de Turing, pois $Y_C g
\equal[\beta] g(Y_C g)$, mas não $Y_C g \bred g(Y_C g)$. Seja $V$ o termo
$\lambd[x][g(xx)]$, de modo que $Y_C \ident \lambd[g][VV]$. Então
\begin{align*}
  VV & \ident (\lambd[x][g(xx)])V \red g(VV)
  \text{ e, portanto,}\\
  Y_C g & \ident (\lambd[g][VV])g \red VV \red g(VV), \text{ mas também}\\
  g(Y_C g) & \ident g((\lambd[g][VV])g) \red g(VV).
\end{align*}
Em outras palavras, $Y_C g$ e $g(Y_C g)$ se reduzem a um termo comum
$g(VV)$; logo, $Y_C g \equal[\beta] g(Y_C g)$. Isso costuma bastar para
as aplicações.

\end{document}
